%======================================================================================
\section{How $H$ and $H_\mathrm{loc}$ are computed}\label{app:hloc}

$\mathrm{MSD}(\tau)\propto\tau^{2H}$, so $\log_{10}\mathrm{MSD}=2H\log_{10}\tau+c$. \textbf{$H$} is half the slope of the
least-squares line over all 21 lags (equal weight per lag). \textbf{$H_\mathrm{loc}(\tau_c)$} is half the slope of the line fitted on
the lags with $|\log_{10}\tau-\log_{10}\tau_c|\le0.25$ (a factor 1.78 on each side), equal weight; where fewer than
3 lags lie in the window, the 3 nearest are used. It is a windowed regression over about half a decade.
$\Delta H$ and $\Delta H_\mathrm{loc}$ of two curves are the differences of the fitted values; for a stratified
average they are the same regression on the averaged $\Delta\log\mathrm{MSD}$ (linear: slope of a difference = difference of slopes).
Bands: the same computation on each bootstrap draw, percentiles lag by lag.

\begin{table}[H]
\centering\footnotesize
\begin{tabular}{@{}r l@{\hspace{14pt}}r l@{\hspace{14pt}}r l@{}}
\toprule
$\tau$ (s) & lags used & $\tau$ (s) & lags used & $\tau$ (s) & lags used \\
\midrule
10 & 10--30 (3 nearest) & 90 & 60--140 & 320 & 210--470 \\
20 & 20--40 (3 nearest) & 100 & 60--170 & 390 & 260--580 \\
30 & 20--50 & 110 & 80--170 & 470 & 320--710 \\
40 & 30--60 & 140 & 80--210 & 580 & 390--1000 \\
50 & 30--80 & 170 & 100--260 & 710 & 470--1000 \\
60 & 40--100 & 210 & 140--320 & 870 & 580--1000 \\
80 & 50--140 & 260 & 170--390 & 1000 & 580--1000 \\
\bottomrule
\end{tabular}
\caption{Lags of the window of each $H_\mathrm{loc}$ value. At 10 and 20~s the window is not centred; at 870 and 1000~s
it is cut at 1000~s and the two values use the same lags, so the end of the curve flattens by construction. The window
resolves nothing narrower than about half a decade.}
\end{table}

%======================================================================================
\section{Which statistics are drawn}\label{app:redundancy}

\textbf{Ratios.} Ratios are sums of numerators over sums of denominators over the runs of a population (a run
goes to the stretch that holds its middle). In a small stratum a glide in rising air has almost no height lost, the glide ratio
diverges and its average over strata is unstable; stratified contrasts use the \emph{height lost per km of glide}, the inverse. Pooled tables keep the glide ratio.

\textbf{Redundancy.} Unit: a population on a site-day with at least 25 flights (119 units). A statistic is redundant when the others explain most of it across units ($R^2$ of its linear
regression on them); backward elimination drops the most explained until every $R^2\le0.5$. The four shares of time sum to 100\%, so one of them goes first.

\begin{table}[H]
\centering\footnotesize
\begin{tabular}{@{}l l c@{}}
\toprule
statistic & status & R$^2$ on the five drawn \\
\midrule
time in climb (\%) & drawn & 0.23 \\
searches that end in a climb (\%) & drawn & 0.41 \\
climb rate (m/s) & drawn & 0.20 \\
glide speed (km/h) & drawn & 0.26 \\
glide ratio & drawn & 0.38 \\
time in transition (\%) & dropped, number 1, R$^2$ 1.00 when dropped & 0.83 \\
height gained per climb (m) & dropped, number 2, R$^2$ 0.98 when dropped & 0.91 \\
climbs that start after a search (\%) & dropped, number 3, R$^2$ 0.93 when dropped & 0.77 \\
cross-country speed over 30 min (km/h) & dropped, number 4, R$^2$ 0.93 when dropped & 0.90 \\
time in a search in vain (\%) & dropped, number 5, R$^2$ 0.88 when dropped & 0.71 \\
climbs per hour & dropped, number 6, R$^2$ 0.72 when dropped & 0.71 \\
time in a search that ends in a climb (\%) & dropped, number 7, R$^2$ 0.65 when dropped & 0.65 \\
\bottomrule
\end{tabular}

\caption{Five statistics remain; the others are explained by them.}
\end{table}

%======================================================================================
\section{Robustness to the threshold}\label{app:robustness}

With $q=0.2$ a contact needs only 26.4~min: groups are more and shorter, and up to one contact in five may be chance.

\begin{table}[H]
\centering\footnotesize
\begin{tabular}{@{}l l l@{}}
\toprule
 & $q=0.1$, $\tau^*=38.1$~min & $q=0.2$, $\tau^*=26.4$~min \\
\midrule
group flights (TA-MSD) & 33\,435 & 48\,777 \\
$H_\mathrm{loc}(1000~\mathrm{s})$, $\Sl$ & 0.837 & 0.836 \\
$H_\mathrm{loc}(1000~\mathrm{s})$, $\A$ & 0.860 & 0.854 \\
$H_\mathrm{loc}(1000~\mathrm{s})$, $\G$ & 0.893 & 0.892 \\
$\Delta H_\mathrm{loc}(1000~\mathrm{s})$, $\G-\A$ same flight & $+$0.037 [$+$0.035, $+$0.039] & $+$0.032 [$+$0.029, $+$0.035] \\
$\Delta H_\mathrm{loc}(1000~\mathrm{s})$, $\G-\A$ same s-d-h & $+$0.043 [$+$0.040, $+$0.046] & $+$0.042 [$+$0.040, $+$0.044] \\
$\Delta H_\mathrm{loc}(1000~\mathrm{s})$, $\A-\Sl$ same s-d-h & $+$0.003 [$+$0.001, $+$0.006] & $+$0.002 [$-$0.001, $+$0.005] \\
$\Delta H_\mathrm{loc}(100~\mathrm{s})$, $\G-\A$ same flight & $-$0.009 [$-$0.009, $-$0.008] & $-$0.016 [$-$0.017, $-$0.015] \\
$\Delta H_\mathrm{loc}(100~\mathrm{s})$, $\G-\A$ same s-d-h & $+$0.002 [$+$0.001, $+$0.002] & $-$0.004 [$-$0.005, $-$0.003] \\
$\Delta$ searches to climb, $\G-\A$ same s-d-h (points) & $+$2.6 [$+$2.3, $+$3.0] & $+$3.4 [$+$3.0, $+$3.7] \\
$\Delta$ climb rate, $\G-\A$ same s-d-h (m/s) & $+$0.029 [$+$0.025, $+$0.034] & $+$0.024 [$+$0.020, $+$0.029] \\
$\Delta$ height lost per km, $\G-\A$ same s-d-h (m/km) & $-$9.5 [$-$10.5, $-$8.6] & $-$11.5 [$-$12.2, $-$10.9] \\
$\Delta$ cross-country speed, $\G-\A$ same s-d-h (km/h) & $+$1.29 [$+$1.16, $+$1.41] & $+$1.14 [$+$1.01, $+$1.30] \\
\bottomrule
\end{tabular}

\caption{The main results with the two thresholds (s-d-h: same site-day and hour). Every result at 1000~s and every
phase statistic keeps its sign; at 100~s the estimates turn slightly negative, consistent with no effect within one phase.}
\end{table}

%======================================================================================
\section{Pooled curves, hang gliders and sizes}\label{app:curves}

\begin{figure}[H]
\centering
\includegraphics[width=0.78\textwidth]{tamsd.pdf}
\caption{TA-MSD and $H_\mathrm{loc}$ of $\Sl$, $\G$ and $\A$, paragliders and hang gliders (hang gliders: 418 group flights, wide bands at long lags).}
\end{figure}

\begin{figure}[H]
\centering
\includegraphics[width=0.78\textwidth]{size_tamsd.pdf}
\caption{By the size of the configuration (orange, darker with more pilots), $\Sl$ (blue) and, dashed, the same flights alone for each size.}
\end{figure}

%======================================================================================
\section{Values behind the figures}\label{app:tables}

The tables and figures that the summary figures of the main text replace, in the order of the text.

\begin{table}[H]
\centering\footnotesize
\begin{tabular}{@{}l l l l@{}}
\toprule
 & solo flights $\Sl$ & group time $\G$ & same flights alone $\A$ \\
\midrule
flights (TA-MSD) & 119\,129 & 33\,435 & 29\,611 \\
$H$ (10--1000~s) & 0.899 [0.898, 0.899] & 0.916 [0.915, 0.917] & 0.914 [0.913, 0.915] \\
$H_\mathrm{loc}(100~\mathrm{s})$ & 0.919 [0.919, 0.919] & 0.927 [0.927, 0.928] & 0.932 [0.931, 0.932] \\
$H_\mathrm{loc}(1000~\mathrm{s})$ & 0.837 [0.837, 0.838] & 0.893 [0.891, 0.894] & 0.860 [0.859, 0.862] \\
time in climb (\%) & 30.8 [30.7, 30.9] & 33.0 [32.8, 33.2] & 29.1 [28.9, 29.2] \\
searches ending in a climb (\%) & 43.3 [43.2, 43.4] & 49.5 [49.2, 49.7] & 45.9 [45.7, 46.1] \\
climb rate (m/s) & 1.22 [1.22, 1.22] & 1.26 [1.25, 1.27] & 1.33 [1.32, 1.34] \\
glide speed (km/h) & 30.1 [30.1, 30.2] & 33.3 [33.1, 33.5] & 32.8 [32.6, 32.9] \\
glide ratio & 10.4 [10.4, 10.4] & 12.9 [12.8, 13.0] & 11.3 [11.2, 11.4] \\
cross-country speed (km/h) & 17.2 [17.2, 17.3] & 20.0 [19.8, 20.3] & 20.0 [19.8, 20.1] \\
\bottomrule
\end{tabular}

\caption{Pooled statistics of $\Sl$, $\G$ and $\A$ (paragliders), with 5--95\% bands (Figure~\ref{fig:raw}). Exponents on
loggers of at most 10~s, phases on 1~Hz flights.}
\label{tab:raw}
\end{table}

\begin{table}[H]
\centering\footnotesize\setlength{\tabcolsep}{3.5pt}
\begin{tabular}{@{}l c c c c c@{}}
\toprule
$\Delta$ & \makecell{$\G-\Sl$\\pooled} & \makecell{$\G-\Sl$\\same s-d-h} & \makecell{$\A-\Sl$\\same s-d-h} & \makecell{$\G-\A$\\same flight} & \makecell{$\G-\A$\\same s-d-h} \\
\midrule
$H_\mathrm{loc}(100~\mathrm{s})$ & \makecell{$+$0.008\\{\scriptsize[$+$0.007, $+$0.009]}} & \makecell{$+$0.002\\{\scriptsize[$+$0.001, $+$0.003]}} & \makecell{$+$0.005\\{\scriptsize[$+$0.004, $+$0.005]}} & \makecell{$-$0.009\\{\scriptsize[$-$0.009, $-$0.008]}} & \makecell{$+$0.002\\{\scriptsize[$+$0.001, $+$0.002]}} \\[2pt]
$H_\mathrm{loc}(1000~\mathrm{s})$ & \makecell{$+$0.055\\{\scriptsize[$+$0.053, $+$0.057]}} & \makecell{$+$0.041\\{\scriptsize[$+$0.038, $+$0.044]}} & \makecell{$+$0.003\\{\scriptsize[$+$0.001, $+$0.006]}} & \makecell{$+$0.037\\{\scriptsize[$+$0.035, $+$0.039]}} & \makecell{$+$0.043\\{\scriptsize[$+$0.040, $+$0.046]}} \\[2pt]
time in climb (\%) & \makecell{$+$2.20\\{\scriptsize[$+$1.96, $+$2.45]}} & \makecell{$+$3.43\\{\scriptsize[$+$3.19, $+$3.63]}} & \makecell{$+$0.43\\{\scriptsize[$+$0.25, $+$0.59]}} & \makecell{$+$5.71\\{\scriptsize[$+$5.52, $+$5.92]}} & \makecell{$+$3.28\\{\scriptsize[$+$3.11, $+$3.47]}} \\[2pt]
searches to climb (\%) & \makecell{$+$6.18\\{\scriptsize[$+$5.82, $+$6.46]}} & \makecell{$+$4.12\\{\scriptsize[$+$3.68, $+$4.54]}} & \makecell{$+$1.60\\{\scriptsize[$+$1.23, $+$1.97]}} & \makecell{$+$8.75\\{\scriptsize[$+$8.37, $+$9.15]}} & \makecell{$+$2.59\\{\scriptsize[$+$2.25, $+$2.99]}} \\[2pt]
climb rate (m/s) & \makecell{$+$0.040\\{\scriptsize[$+$0.030, $+$0.051]}} & \makecell{$+$0.031\\{\scriptsize[$+$0.025, $+$0.039]}} & \makecell{$+$0.007\\{\scriptsize[$+$0.001, $+$0.012]}} & \makecell{$+$0.024\\{\scriptsize[$+$0.019, $+$0.031]}} & \makecell{$+$0.029\\{\scriptsize[$+$0.025, $+$0.034]}} \\[2pt]
glide speed (km/h) & \makecell{$+$3.18\\{\scriptsize[$+$2.97, $+$3.39]}} & \makecell{$+$1.97\\{\scriptsize[$+$1.82, $+$2.10]}} & \makecell{$+$1.15\\{\scriptsize[$+$1.05, $+$1.27]}} & \makecell{$+$1.36\\{\scriptsize[$+$1.27, $+$1.46]}} & \makecell{$+$0.87\\{\scriptsize[$+$0.79, $+$0.98]}} \\[2pt]
height lost per km (m/km) & \makecell{$-$18.6\\{\scriptsize[$-$19.3, $-$17.8]}} & \makecell{$-$10.0\\{\scriptsize[$-$10.7, $-$9.1]}} & \makecell{$-$0.5\\{\scriptsize[$-$1.1, $+$0.3]}} & \makecell{$-$21.1\\{\scriptsize[$-$22.1, $-$20.1]}} & \makecell{$-$9.5\\{\scriptsize[$-$10.5, $-$8.6]}} \\[2pt]
cross-country speed (km/h) & \makecell{$+$2.81\\{\scriptsize[$+$2.59, $+$3.01]}} & \makecell{$+$2.25\\{\scriptsize[$+$2.09, $+$2.39]}} & \makecell{$+$1.31\\{\scriptsize[$+$1.21, $+$1.43]}} & \makecell{$+$0.70\\{\scriptsize[$+$0.60, $+$0.78]}} & \makecell{$+$1.29\\{\scriptsize[$+$1.16, $+$1.41]}} \\[2pt]
\bottomrule
\end{tabular}

\caption{The five contrasts of Section~\ref{sec:other} (Figure~\ref{fig:decompose}): difference and 5--95\% band
(s-d-h: same site-day and hour).}
\label{tab:decompose}
\end{table}

\begin{table}[H]
\centering\footnotesize\setlength{\tabcolsep}{2.6pt}
\begin{tabular}{@{}l r c c c c c c c c c@{}}
\toprule
 & & \multicolumn{3}{c}{fitted $H$ (10--1000 s)} & \multicolumn{3}{c}{$\Delta H$} & \multicolumn{3}{c}{$\Delta H_\mathrm{loc}(1000~\mathrm{s})$} \\
\cmidrule(lr){3-5}\cmidrule(lr){6-8}\cmidrule(l){9-11}
compared within & \makecell{G\\kept (\%)} & $\G$ & $\Sl\mid\G$ & $\A\mid\G$ & \makecell{$\G-\Sl$\\as is} & \makecell{$\G-\Sl$\\reweighted} & \makecell{$\G-\A$\\reweighted} & \makecell{$\G-\Sl$\\as is} & \makecell{$\G-\Sl$\\reweighted} & \makecell{$\G-\A$\\reweighted} \\
\midrule
none & 100 & 0.916 & 0.899 & 0.914 & $+$0.017 & \makecell{$+$0.017\\{\scriptsize[$+$0.016, $+$0.018]}} & \makecell{$+$0.002\\{\scriptsize[$+$0.001, $+$0.003]}} & $+$0.055 & \makecell{$+$0.055\\{\scriptsize[$+$0.053, $+$0.057]}} & \makecell{$+$0.032\\{\scriptsize[$+$0.030, $+$0.034]}} \\[2pt]
zone & 100 & 0.919 & 0.902 & 0.916 & $+$0.021 & \makecell{$+$0.018\\{\scriptsize[$+$0.017, $+$0.019]}} & \makecell{$+$0.003\\{\scriptsize[$+$0.003, $+$0.004]}} & $+$0.059 & \makecell{$+$0.056\\{\scriptsize[$+$0.054, $+$0.058]}} & \makecell{$+$0.028\\{\scriptsize[$+$0.026, $+$0.029]}} \\[2pt]
zone $\times$ month & 99 & 0.919 & 0.902 & 0.915 & $+$0.022 & \makecell{$+$0.017\\{\scriptsize[$+$0.016, $+$0.018]}} & \makecell{$+$0.003\\{\scriptsize[$+$0.003, $+$0.004]}} & $+$0.064 & \makecell{$+$0.054\\{\scriptsize[$+$0.053, $+$0.057]}} & \makecell{$+$0.028\\{\scriptsize[$+$0.027, $+$0.031]}} \\[2pt]
zone $\times$ hour & 99 & 0.919 & 0.902 & 0.913 & $+$0.022 & \makecell{$+$0.017\\{\scriptsize[$+$0.016, $+$0.018]}} & \makecell{$+$0.005\\{\scriptsize[$+$0.005, $+$0.006]}} & $+$0.064 & \makecell{$+$0.052\\{\scriptsize[$+$0.050, $+$0.054]}} & \makecell{$+$0.026\\{\scriptsize[$+$0.025, $+$0.028]}} \\[2pt]
site-day $\times$ hour & 53 & 0.917 & 0.904 & 0.913 & $+$0.015 & \makecell{$+$0.013\\{\scriptsize[$+$0.012, $+$0.014]}} & \makecell{$+$0.008\\{\scriptsize[$+$0.007, $+$0.009]}} & $+$0.051 & \makecell{$+$0.044\\{\scriptsize[$+$0.041, $+$0.046]}} & \makecell{$+$0.033\\{\scriptsize[$+$0.031, $+$0.036]}} \\[2pt]
\bottomrule
\end{tabular}

\caption{Fitted $H$ and differences by stratification (Figure~\ref{fig:zones-forest}). ``Reweighted'': the other population
reweighted to the strata of $\G$ (value, 5--95\% band); ``as is'': the same strata without reweighting. $H_{\Sl\mid\G}$: $H$
of $\Sl$ reweighted to $\G$. The columns $\G$ and ``kept'' refer to the strata shared with $\Sl$; on those shared with $\A$
they differ slightly.}
\label{tab:zones}
\end{table}

\begin{figure}[H]
\centering
\includegraphics[width=\textwidth]{zones_curves.pdf}
\caption{$H_\mathrm{loc}(\tau)$ of $\G$ (orange), of $\Sl$ reweighted to the strata of $\G$ (blue) and of $\A$ reweighted to
them (green), 5--95\% bands (Section~\ref{sec:zones}). First panel: all flights pooled. ``Share of G kept'': flights of $\G$
with a solo counterpart in their stratum.}
\label{fig:zones}
\end{figure}

\begin{figure}[H]
\centering
\includegraphics[width=\textwidth]{report_route.pdf}
\caption{$\Delta\Hloc(1000~\mathrm{s})$ of the matched pieces against the solo flights of their site-day
and hour (blue) and against the same route flown alone (purple; open: same season and hour), by sample,
radius $\rho$ and cover. First block: the sample of Section~\ref{sec:route}.}
\label{fig:route-forest}
\end{figure}

\begin{table}[H]
\centering\footnotesize\setlength{\tabcolsep}{3pt}
\begin{tabular}{@{}l r c c c@{}}
\toprule
$\rho$, cover & pieces & route, any date & route, same season and hour & site-day and hour \\
\midrule
500~m, 80\% & 10\,000: 4\,933 & $+$0.018 [$+$0.014, $+$0.021] & $+$0.016 [$+$0.012, $+$0.021] & $+$0.060 [$+$0.055, $+$0.066] \\
500~m, 80\% & 2\,000: 958 & $+$0.021 [$+$0.013, $+$0.027] & $+$0.017 [$+$0.008, $+$0.024] & $+$0.061 [$+$0.051, $+$0.072] \\
300~m, 80\% & 2\,000: 675 & $+$0.022 [$+$0.012, $+$0.030] & $+$0.018 [$+$0.008, $+$0.029] & $+$0.066 [$+$0.053, $+$0.079] \\
1000~m, 80\% & 2\,000: 1\,284 & $+$0.014 [$+$0.010, $+$0.021] & $+$0.016 [$+$0.010, $+$0.023] & $+$0.057 [$+$0.049, $+$0.065] \\
500~m, 70\% & 2\,000: 1\,062 & $+$0.017 [$+$0.011, $+$0.025] & $+$0.016 [$+$0.009, $+$0.024] & $+$0.061 [$+$0.051, $+$0.070] \\
500~m, 90\% & 2\,000: 793 & $+$0.020 [$+$0.012, $+$0.028] & $+$0.015 [$+$0.006, $+$0.024] & $+$0.062 [$+$0.050, $+$0.073] \\
\bottomrule
\end{tabular}

\caption{$\Delta H_\mathrm{loc}(1000~\mathrm{s})$ of the matched pieces against the same route flown alone and against the
solo flights of their site-day and hour, by radius $\rho$ and cover (Figure~\ref{fig:route-forest}). Pieces: drawn at random,
and matched.}
\label{tab:route}
\end{table}

\begin{table}[H]
\centering\footnotesize\setlength{\tabcolsep}{4pt}
\begin{tabular}{@{}l r c c c@{}}
\toprule
comparison & events & $\Delta H$ & $\Delta H_\mathrm{loc}(100)$ & $\Delta H_\mathrm{loc}(1000)$ \\
\midrule
entering: start $-$ before & 6\,524 & \makecell{$+$0.010\\{\scriptsize[$+$0.009, $+$0.011]}} & \makecell{$+$0.005\\{\scriptsize[$+$0.004, $+$0.006]}} & \makecell{$+$0.027\\{\scriptsize[$+$0.023, $+$0.031]}} \\[2pt]
leaving: end $-$ after & 13\,643 & \makecell{$+$0.005\\{\scriptsize[$+$0.004, $+$0.006]}} & \makecell{$-$0.001\\{\scriptsize[$-$0.002, $-$0.000]}} & \makecell{$+$0.031\\{\scriptsize[$+$0.028, $+$0.033]}} \\[2pt]
after $-$ before & 6\,320 & \makecell{$+$0.004\\{\scriptsize[$+$0.003, $+$0.005]}} & \makecell{$+$0.006\\{\scriptsize[$+$0.005, $+$0.007]}} & \makecell{$-$0.003\\{\scriptsize[$-$0.007, $+$0.001]}} \\[2pt]
control A: later $-$ earlier & 16\,998 & \makecell{$+$0.005\\{\scriptsize[$+$0.004, $+$0.005]}} & \makecell{$+$0.005\\{\scriptsize[$+$0.005, $+$0.006]}} & \makecell{$-$0.002\\{\scriptsize[$-$0.005, $+$0.000]}} \\[2pt]
control S: later $-$ earlier & 69\,648 & \makecell{$+$0.009\\{\scriptsize[$+$0.008, $+$0.009]}} & \makecell{$+$0.010\\{\scriptsize[$+$0.010, $+$0.011]}} & \makecell{$-$0.002\\{\scriptsize[$-$0.003, $-$0.000]}} \\[2pt]
\bottomrule
\end{tabular}

\caption{Differences of the fitted $H$ and of $H_\mathrm{loc}$ at 100 and 1000~s between windows, on the events with both
windows (Figure~\ref{fig:prepost-forest}). Controls: later minus earlier; leaving is end minus after, so its drift is the
control with the opposite sign.}
\label{tab:prepost}
\end{table}

\begin{table}[H]
\centering\footnotesize\setlength{\tabcolsep}{3.5pt}
\begin{tabular}{@{}l r c c c c c@{}}
\toprule
comparison & events & \makecell{time in\\climb (points)} & \makecell{searches to\\climb (points)} & \makecell{climb rate\\(m/s)} & \makecell{glide speed\\(km/h)} & \makecell{height lost\\per km (m/km)} \\
\midrule
entering: start $-$ before & 4\,786 & \makecell{$+$0.7\\{\scriptsize[$+$0.3, $+$1.1]}} & \makecell{$+$0.0\\{\scriptsize[$-$1.1, $+$1.0]}} & \makecell{$+$0.046\\{\scriptsize[$+$0.034, $+$0.062]}} & \makecell{$+$0.92\\{\scriptsize[$+$0.69, $+$1.15]}} & \makecell{$-$2.3\\{\scriptsize[$-$3.5, $-$1.3]}} \\[2pt]
leaving: end $-$ after & 10\,097 & \makecell{$+$1.5\\{\scriptsize[$+$1.2, $+$1.9]}} & \makecell{$+$1.2\\{\scriptsize[$+$0.4, $+$1.9]}} & \makecell{$-$0.007\\{\scriptsize[$-$0.017, $+$0.003]}} & \makecell{$+$0.52\\{\scriptsize[$+$0.35, $+$0.64]}} & \makecell{$-$10.9\\{\scriptsize[$-$11.8, $-$10.0]}} \\[2pt]
after $-$ before & 4\,718 & \makecell{$-$1.4\\{\scriptsize[$-$1.8, $-$1.0]}} & \makecell{$-$2.8\\{\scriptsize[$-$4.0, $-$1.6]}} & \makecell{$+$0.043\\{\scriptsize[$+$0.024, $+$0.059]}} & \makecell{$+$0.33\\{\scriptsize[$+$0.12, $+$0.55]}} & \makecell{$+$9.1\\{\scriptsize[$+$7.8, $+$10.4]}} \\[2pt]
control A: later $-$ earlier & 12\,568 & \makecell{$-$1.1\\{\scriptsize[$-$1.3, $-$0.9]}} & \makecell{$-$1.4\\{\scriptsize[$-$2.0, $-$0.8]}} & \makecell{$+$0.018\\{\scriptsize[$+$0.009, $+$0.027]}} & \makecell{$+$0.07\\{\scriptsize[$-$0.05, $+$0.20]}} & \makecell{$+$3.3\\{\scriptsize[$+$2.5, $+$4.2]}} \\[2pt]
control S: later $-$ earlier & 49\,092 & \makecell{$-$1.6\\{\scriptsize[$-$1.7, $-$1.5]}} & \makecell{$-$2.3\\{\scriptsize[$-$2.7, $-$2.0]}} & \makecell{$+$0.011\\{\scriptsize[$+$0.008, $+$0.015]}} & \makecell{$+$0.04\\{\scriptsize[$-$0.03, $+$0.10]}} & \makecell{$+$4.9\\{\scriptsize[$+$4.5, $+$5.3]}} \\[2pt]
\bottomrule
\end{tabular}

\caption{The phase statistics of the same comparisons (differences with 5--95\% band); runs and glides are placed in the
window that holds their middle.}
\label{tab:prepost-phase}
\end{table}

\begin{figure}[H]
\centering
\includegraphics[width=0.86\textwidth]{prepost_3000_profile.pdf}
\caption{The four windows of Section~\ref{sec:prepost} on the events that have all of them. Start and end overlap when the
group lasts less than $2W$.}
\label{fig:prepost-profile}
\end{figure}

\begin{table}[H]
\centering\footnotesize\setlength{\tabcolsep}{3pt}
\begin{tabular}{@{}l r c c c c c c c c c@{}}
\toprule
 & & \multicolumn{2}{c}{$H$} & \multicolumn{2}{c}{$H_\mathrm{loc}(1000~\mathrm{s})$} & & & & & \\
\cmidrule(lr){3-4}\cmidrule(lr){5-6}
pilots & flights & $\G_n$ & alone & $\G_n$ & alone & \makecell{time in\\climb (\%)} & \makecell{searches\\to climb (\%)} & \makecell{climb\\rate (m/s)} & \makecell{glide speed\\(km/h)} & \makecell{glide\\ratio} \\
\midrule
solo & 119\,129 & 0.899 &  & 0.837 &  & 30.8 & 43.3 & 1.22 & 30.1 & 10.4 \\
2 & 26\,411 & 0.917 & 0.914 & 0.889 & 0.861 & 32.7 & 49.0 & 1.27 & 33.2 & 12.5 \\
3 & 9\,480 & 0.923 & 0.919 & 0.906 & 0.867 & 33.6 & 49.9 & 1.26 & 33.5 & 13.1 \\
4 & 4\,073 & 0.924 & 0.920 & 0.910 & 0.865 & 33.6 & 51.0 & 1.25 & 33.5 & 13.7 \\
5--6 & 2\,686 & 0.925 & 0.922 & 0.915 & 0.870 & 33.1 & 50.1 & 1.23 & 33.1 & 14.1 \\
7+ & 1\,142 & 0.933 & 0.924 & 0.944 & 0.871 & 34.6 & 51.7 & 1.20 & 34.6 & 15.9 \\
\bottomrule
\end{tabular}

\caption{Paragliders, pooled by the size of the configuration (``alone'': the same flights alone).}
\label{tab:size}
\end{table}

\begin{table}[H]
\centering\footnotesize
\resizebox{\textwidth}{!}{\begin{tabular}{@{}l *{9}{l}@{}}
\toprule
 & \multicolumn{5}{c}{$\G_n-\A$, same flight} & \multicolumn{4}{c}{$\G_n-\G_2$, same site-day-hour} \\
\cmidrule(lr){2-6}\cmidrule(l){7-10}
$n$ & 2 & 3 & 4 & 5--6 & 7+ & 3 & 4 & 5--6 & 7+ \\
\midrule
$H_\mathrm{loc}(100~\mathrm{s})$ & $-$0.008 & $-$0.014 & $-$0.012 & $-$0.021 & $-$0.013 & $-$0.003 & $-$0.002 & $-$0.013 & $+$0.009 \\
$H_\mathrm{loc}(1000~\mathrm{s})$ & $+$0.026 & $+$0.027 & $+$0.043 & $+$0.038 & $+$0.087 & $+$0.004 & $-$0.001 & $-$0.015 & $+$0.064 \\
time in climb (\%) & $+$5.76 & $+$8.26 & $+$9.37 & $+$9.11 & $+$9.91 & $+$1.98 & $+$3.77 & $+$3.54 & $+$4.03 \\
searches ending in a climb (\%) & $+$8.25 & $+$9.41 & $+$10.78 & $+$9.97 & $+$11.25 & $+$0.28 & $+$0.55 & $+$0.35 & $+$1.19 \\
climb rate (m/s) & $+$0.021 & $+$0.006 & $-$0.008 & $-$0.012 & $-$0.055 & $+$0.002 & $-$0.003 & $-$0.006 & $-$0.017 \\
glide speed (km/h) & $+$1.29 & $+$1.15 & $+$1.32 & $+$0.93 & $+$0.96 & $-$0.23 & $-$0.40 & $-$0.46 & $-$0.18 \\
height lost per km of glide (m/km) & $-$17.6 & $-$22.8 & $-$28.1 & $-$32.3 & $-$28.1 & $-$3.0 & $-$8.0 & $-$8.7 & $-$7.5 \\
\bottomrule
\end{tabular}}

\caption{The contrasts by size of Section~\ref{sec:size} (Figure~\ref{fig:size}).}
\label{tab:size-contrasts}
\end{table}

\begin{table}[H]
\centering\footnotesize
\begin{tabular}{@{}l l c c@{}}
\toprule
test & comparison & $\Delta\Hloc(1000~\mathrm{s})$ & band \\
\midrule
raw & $\G-\Sl$ pooled & \DHlocThousandRaw & \DHlocThousandRawBand \\
pilots, day, site, hour & $\G-\Sl$ same site-day and hour & \DHlocThousandSdh & \DHlocThousandSdhBand \\
 & $\A-\Sl$ same site-day and hour (selection) & \DHlocThousandSel & \DHlocThousandSelBand \\
 & $\G-\A$ same flight & \DHlocThousandFlight & \DHlocThousandFlightBand \\
 & $\G-\A$ same site-day and hour & \DHlocThousandGroupSdh & \DHlocThousandGroupSdhBand \\
place & $\G-\Sl$ reweighted, site-day$\times$hour & \ZoneDHlocSSdh & \ZoneDHlocSSdhBand \\
 & $\G-\A$ reweighted, site-day$\times$hour & \ZoneDHlocASdh & \ZoneDHlocASdhBand \\
route & matched pieces $-$ $\Sl$ same site-day and hour & \RouteBaseDHloc & \RouteBaseDHlocBand \\
 & matched pieces $-$ same route alone & \RouteDHloc & \RouteDHlocBand \\
 & the same, same season and hour & \RouteDHlocSeason & \RouteDHlocSeasonBand \\
moment & entering, start $-$ before & \PPEntryDHloc & \PPEntryDHlocBand \\
 & leaving, end $-$ after & \PPExitDHloc & \PPExitDHlocBand \\
 & after $-$ before & \PPAfterBeforeDHloc & \PPAfterBeforeDHlocBand \\
 & drift, solo time of group flights ($\A$) & \PPCtrlADHloc & \PPCtrlADHlocBand \\
 & drift, solo flights ($\Sl$) & \PPCtrlSDHloc & \PPCtrlSDHlocBand \\
\bottomrule
\end{tabular}
\caption{Every estimate of the long-lag difference, from the raw one to the four tests. The route
rows refer to the matched pieces only. The drifts are later minus earlier windows; for leaving they
enter with the opposite sign.}
\label{tab:estimates}
\end{table}

%======================================================================================
\section{Before, during and after: curves and short windows}\label{app:prepost}

\begin{figure}[H]
\centering
\includegraphics[width=\textwidth]{prepost_3000_curves.pdf}
\caption{$\Hloc(\tau)$ of the same flights alone (green) and in the group (orange), entering (left)
and leaving (middle), windows of \PPWindowMin{} min, 5--95\% bands. Right: the differences, group minus alone, with the two
controls (later minus earlier window; dashed); for leaving, the drift to compare is the control with the opposite sign.}
\label{fig:prepost}
\end{figure}

Same test as Section~\ref{sec:prepost} with windows of \PPShortWindowMin{} min, the shortest a piece of the TA-MSD allows.
Groups of at least \PPShortWindowMin{} min enter: more events, but at 1000~s a window holds fewer than thirty origins.
At 1000~s entering then stays just above its drifts and leaving at them, taken with the opposite sign
(\PPShortEntryDHloc{} entering, \PPShortExitDHloc{} leaving, against drifts of \PPShortCtrlADHloc{} and \PPShortCtrlSDHloc);
on the fitted $H$, dominated by short lags, entering stays below its drifts and leaving above them.
The phase statistics at entry depend on the window (time in climb \PPShortEntryClimb{} points, glide speed
\PPShortEntrySpeed{} km/h here, against \PPEntryClimb{} and \PPEntrySpeed{} with \PPWindowMin{} min).

\begin{figure}[H]
\centering
\includegraphics[width=\textwidth]{prepost_1260_curves.pdf}
\caption{As Figure~\ref{fig:prepost}, windows of \PPShortWindowMin{} min.}
\end{figure}

\begin{table}[H]
\centering\footnotesize\setlength{\tabcolsep}{4pt}
\begin{tabular}{@{}l r c c c@{}}
\toprule
comparison & events & $\Delta H$ & $\Delta H_\mathrm{loc}(100)$ & $\Delta H_\mathrm{loc}(1000)$ \\
\midrule
entering: start $-$ before & 19\,876 & \makecell{$+$0.002\\{\scriptsize[$+$0.001, $+$0.004]}} & \makecell{$-$0.002\\{\scriptsize[$-$0.004, $-$0.001]}} & \makecell{$-$0.003\\{\scriptsize[$-$0.007, $+$0.000]}} \\[2pt]
leaving: end $-$ after & 30\,884 & \makecell{$+$0.002\\{\scriptsize[$+$0.001, $+$0.003]}} & \makecell{$-$0.005\\{\scriptsize[$-$0.006, $-$0.004]}} & \makecell{$+$0.008\\{\scriptsize[$+$0.006, $+$0.011]}} \\[2pt]
after $-$ before & 14\,482 & \makecell{$+$0.004\\{\scriptsize[$+$0.003, $+$0.005]}} & \makecell{$+$0.006\\{\scriptsize[$+$0.005, $+$0.008]}} & \makecell{$+$0.002\\{\scriptsize[$-$0.001, $+$0.006]}} \\[2pt]
control A: later $-$ earlier & 33\,658 & \makecell{$+$0.006\\{\scriptsize[$+$0.005, $+$0.006]}} & \makecell{$+$0.006\\{\scriptsize[$+$0.006, $+$0.007]}} & \makecell{$-$0.011\\{\scriptsize[$-$0.015, $-$0.008]}} \\[2pt]
control S: later $-$ earlier & 107\,324 & \makecell{$+$0.009\\{\scriptsize[$+$0.009, $+$0.010]}} & \makecell{$+$0.010\\{\scriptsize[$+$0.010, $+$0.011]}} & \makecell{$-$0.007\\{\scriptsize[$-$0.008, $-$0.005]}} \\[2pt]
\bottomrule
\end{tabular}

\caption{As Table~\ref{tab:prepost}, windows of \PPShortWindowMin{} min.}
\end{table}

\begin{table}[H]
\centering\footnotesize\setlength{\tabcolsep}{3.5pt}
\begin{tabular}{@{}l r c c c c c@{}}
\toprule
comparison & events & \makecell{time in\\climb (points)} & \makecell{searches to\\climb (points)} & \makecell{climb rate\\(m/s)} & \makecell{glide speed\\(km/h)} & \makecell{height lost\\per km (m/km)} \\
\midrule
entering: start $-$ before & 14\,798 & \makecell{$+$3.3\\{\scriptsize[$+$2.9, $+$3.8]}} & \makecell{$+$1.2\\{\scriptsize[$+$0.3, $+$2.2]}} & \makecell{$+$0.043\\{\scriptsize[$+$0.035, $+$0.053]}} & \makecell{$+$0.63\\{\scriptsize[$+$0.45, $+$0.77]}} & \makecell{$-$8.9\\{\scriptsize[$-$12.1, $-$6.2]}} \\[2pt]
leaving: end $-$ after & 23\,014 & \makecell{$+$2.7\\{\scriptsize[$+$2.3, $+$3.1]}} & \makecell{$+$0.9\\{\scriptsize[$+$0.2, $+$1.6]}} & \makecell{$+$0.008\\{\scriptsize[$+$0.000, $+$0.018]}} & \makecell{$+$0.54\\{\scriptsize[$+$0.42, $+$0.67]}} & \makecell{$-$13.1\\{\scriptsize[$-$14.5, $-$11.5]}} \\[2pt]
after $-$ before & 10\,898 & \makecell{$-$1.5\\{\scriptsize[$-$2.0, $-$1.1]}} & \makecell{$-$4.1\\{\scriptsize[$-$5.2, $-$2.9]}} & \makecell{$+$0.017\\{\scriptsize[$-$0.001, $+$0.030]}} & \makecell{$+$0.17\\{\scriptsize[$-$0.01, $+$0.35]}} & \makecell{$+$9.6\\{\scriptsize[$+$5.7, $+$12.9]}} \\[2pt]
control A: later $-$ earlier & 25\,121 & \makecell{$-$2.1\\{\scriptsize[$-$2.4, $-$1.8]}} & \makecell{$-$2.8\\{\scriptsize[$-$3.5, $-$2.0]}} & \makecell{$-$0.001\\{\scriptsize[$-$0.009, $+$0.006]}} & \makecell{$-$0.29\\{\scriptsize[$-$0.40, $-$0.15]}} & \makecell{$+$3.4\\{\scriptsize[$+$1.9, $+$5.2]}} \\[2pt]
control S: later $-$ earlier & 75\,984 & \makecell{$-$2.0\\{\scriptsize[$-$2.2, $-$1.8]}} & \makecell{$-$3.0\\{\scriptsize[$-$3.5, $-$2.7]}} & \makecell{$-$0.012\\{\scriptsize[$-$0.017, $-$0.009]}} & \makecell{$-$0.31\\{\scriptsize[$-$0.37, $-$0.24]}} & \makecell{$+$6.2\\{\scriptsize[$+$5.0, $+$7.5]}} \\[2pt]
\bottomrule
\end{tabular}

\caption{As Table~\ref{tab:prepost-phase}, windows of \PPShortWindowMin{} min.}
\end{table}

%======================================================================================
\section{Single site-days}\label{app:sitedays}

Six site-days chosen by a fixed rule (most sizes with 15 flights or more, then most flights, at most two per takeoff site). Bands resample flights (too narrow: flights of one group are not independent).

\begin{figure}[H]
\centering
\includegraphics[width=0.9\textwidth]{siteday_tamsd.pdf}
\caption{TA-MSD and $H_\mathrm{loc}$ on six site-days by the size of the configuration, with the solo flights of the same site-day where they reach 15 flights.}
\end{figure}

\begin{table}[H]
\centering\footnotesize
\begin{tabular}{@{}l l l l l l l@{}}
\toprule
site-day & solo & 2 & 3 & 4 & 5--6 & 7+ \\
\midrule
Saint Hilaire sud, 2022-03-26 & 0.750 & 0.853 & 0.894 & 0.818 &  & 0.994 \\
Saint Hilaire sud, 2017-03-28 & 0.854 & 0.931 & 0.901 & 0.956 &  & 0.873 \\
M\'eruz, 2019-05-01 &  & 0.926 & 0.939 &  & 0.938 &  \\
M\'eruz, 2016-04-28 &  & 0.894 & 0.958 &  & 0.957 &  \\
Chalvet Ouest, 2016-07-19 &  & 0.951 &  &  & 0.892 &  \\
La Sambuy, 2020-07-25 &  & 0.941 & 0.901 &  &  &  \\
\bottomrule
\end{tabular}

\caption{$H_\mathrm{loc}(1000~\mathrm{s})$ by size on each site-day (empty: fewer than 15 flights). Differences between days exceed those between sizes (the climb rate goes from 0.79 to 2.11~m/s across the six days).}
\end{table}

%======================================================================================
\section{Leader minus follower}\label{app:leaders}

\begin{figure}[H]
\centering
\includegraphics[width=0.95\textwidth]{mechanisms_leaders.pdf}
\caption{Leader minus follower, same configuration of two pilots.}
\end{figure}

\begin{table}[H]
\centering\footnotesize
\begin{tabular}{@{}l l@{}}
\toprule
$\Delta$ & leader $-$ follower \\
\midrule
$H_\mathrm{loc}(100~\mathrm{s})$ & $-$0.002 [$-$0.002, $-$0.001] \\
$H_\mathrm{loc}(1000~\mathrm{s})$ & $-$0.003 [$-$0.004, $-$0.001] \\
time in climb (\%) & $+$0.91 [$+$0.57, $+$1.25] \\
searches ending in a climb (\%) & $-$0.01 [$-$1.03, $+$0.83] \\
climb rate (m/s) & $+$0.021 [$+$0.014, $+$0.028] \\
glide speed (km/h) & $+$0.86 [$+$0.74, $+$0.99] \\
height lost per km of glide (m/km) & $+$1.6 [$-$0.9, $+$4.5] \\
cross-country speed (km/h) & $+$0.14 [$+$0.11, $+$0.16] \\
\bottomrule
\end{tabular}

\caption{Leader minus follower, same configuration of two pilots (difference with 5--95\% band).}
\end{table}

%======================================================================================
\section{Checks of the viewer and of the analyses}\label{app:audit}

The checks of the grouping itself, its code against the note and the properties of its tables, are in
section 8 of the note.

\begin{itemize}
  \item \textbf{Viewer.} The \emph{Group flights} tab reads \texttt{groups/} of its data folder: its seven tables are \textbf{byte-identical} to those of the chosen run; the tracks it draws (\texttt{SSD\_DATA}) have the same fixes,
        times, positions and trimmed starts as those the grouping read (\texttt{SSD\_DISANTE}).
  \item \textbf{Analyses.} Tests on synthetic flights for every step, and on every number of the README of the analyses; the decomposition is checked on a synthetic case where the day raises
        a statistic by 10 and the group by 2: pooled 10, stratified 2.
\end{itemize}

%======================================================================================
\section{Reproducing the report}\label{app:reproduce}

From \texttt{xc-lab/analysis/groups/group-identification}, SSD\_DISANTE mounted:
\begin{verbatim}
uv run python tamsd.py; uv run python phase_msd.py; uv run python survival.py
uv run python budget.py; uv run python performance.py; uv run python blocks.py
uv run python redundancy.py; uv run python size.py; uv run python decompose.py
uv run python zones.py; uv run python prepost.py; uv run python prepost.py --window 1260
uv run python mechanisms.py; uv run python report/facts.py
uv run python route.py --sample 10000
uv run python report/tables.py; uv run python report/figures.py
cd report; latexmk -pdf group_vs_solo.tex
\end{verbatim}
Heavy pieces stay on SSD\_DISANTE, in the \texttt{analyses/} folder of each run; figures and JSON in \texttt{out/<run>/}.
